% GENERATED FILE. Edit canonical lessons, then run npm run generate:books.
% canonical-source-hash: fnv1a64:e81b91b1eefcbce6
% canonical-lessons: JA-C77-futsuka, JA-C77-mikka, JA-C77-yokka, JA-C77-itsuka, JA-C77-muika

\chapter{Second of the month, Third of the month, Fourth of the month, Fifth of the month, Sixth of the month}
\label{ch:ja-second-of-the-month-third-of-the-month-fourth-of-the-month-fifth-of-the-month-sixth-of-the-month}
\hlchaptermodality{77}

\begin{chapteropening}
\textbf{By the end of this chapter:} \emph{I can say the second of the month, the third of the month, the fourth of the month, the fifth of the month and the sixth of the month.}

Complete the last lesson of chapter 77: I can say the second of the month, the third of the month, the fourth of the month, the fifth of the month and the sixth of the month.
\end{chapteropening}

\section[\texorpdfstring{\ja{ふつか} (futsuka)}{ふつか (futsuka)}]{\ja{ふつか} (futsuka) — the second of the month}

\label{lesson:JA-C77-futsuka}



\begin{quote}
Before the new one: say the Japanese for half a day, then the Japanese for the first of the month.
\end{quote}

\subsection*{You'll want to know: \ja{ふつか}}
\textbf{\ja{ふつか}} — \emph{futsuka} — \textquotedblleft{}the second of the month\textquotedblright{}.

\begin{grammarlens}[title={What you've built}]
One more word for this chapter.
\end{grammarlens}

\subsection*{Your turn}
\begin{itemize}
\raggedright
  \item \emph{Say it:} \emph{futsuka}
  \item \emph{Say it:} \emph{futsuka}, once more
  \item \emph{Say it:} say \emph{futsuka}
  \item \emph{Recall:} say the Japanese for half a day, then the Japanese for the first of the month, then say \emph{futsuka} again
\end{itemize}

\begin{tcolorbox}[breakable,colback=teal!4,colframe=teal!35!black,title={Before you move on}]
What does \ja{ふつか} mean? (\textquotedblleft{}The second of the month\textquotedblright{}.) Say it once more.
\end{tcolorbox}
\section[\texorpdfstring{\ja{みっか} (mikka)}{みっか (mikka)}]{\ja{みっか} (mikka) — the third of the month}

\label{lesson:JA-C77-mikka}



\begin{quote}
Before the new one: say the Japanese for the first of the month, then the Japanese for the second of the month.
\end{quote}

\subsection*{You'll want to know: \ja{みっか}}
\textbf{\ja{みっか}} — \emph{mikka} — \textquotedblleft{}the third of the month\textquotedblright{}.

\begin{grammarlens}[title={What you've built}]
One more word for this chapter.
\end{grammarlens}

\subsection*{Your turn}
\begin{itemize}
\raggedright
  \item \emph{Say it:} \emph{mikka}
  \item \emph{Say it:} \emph{mikka}, once more
  \item \emph{Say it:} say \emph{mikka}
  \item \emph{Recall:} say the Japanese for the first of the month, then the Japanese for the second of the month, then say \emph{mikka} again
\end{itemize}

\begin{tcolorbox}[breakable,colback=teal!4,colframe=teal!35!black,title={Before you move on}]
What does \ja{みっか} mean? (\textquotedblleft{}The third of the month\textquotedblright{}.) Say it once more.
\end{tcolorbox}
\section[\texorpdfstring{\ja{よっか} (yokka)}{よっか (yokka)}]{\ja{よっか} (yokka) — the fourth of the month}

\label{lesson:JA-C77-yokka}



\begin{quote}
Before the new one: say the Japanese for the second of the month, then the Japanese for the third of the month.
\end{quote}

\subsection*{You'll want to know: \ja{よっか}}
\textbf{\ja{よっか}} — \emph{yokka} — \textquotedblleft{}the fourth of the month\textquotedblright{}.

\begin{grammarlens}[title={What you've built}]
One more word for this chapter.
\end{grammarlens}

\subsection*{Your turn}
\begin{itemize}
\raggedright
  \item \emph{Say it:} \emph{yokka}
  \item \emph{Say it:} \emph{yokka}, once more
  \item \emph{Say it:} say \emph{yokka}
  \item \emph{Recall:} say the Japanese for the second of the month, then the Japanese for the third of the month, then say \emph{yokka} again
\end{itemize}

\begin{tcolorbox}[breakable,colback=teal!4,colframe=teal!35!black,title={Before you move on}]
What does \ja{よっか} mean? (\textquotedblleft{}The fourth of the month\textquotedblright{}.) Say it once more.
\end{tcolorbox}
\section[\texorpdfstring{\ja{いつか} (itsuka)}{いつか (itsuka)}]{\ja{いつか} (itsuka) — the fifth of the month; some day}

\label{lesson:JA-C77-itsuka}



\begin{quote}
Before the new one: say the Japanese for the third of the month, then the Japanese for the fourth of the month.
\end{quote}

\subsection*{You'll want to know: \ja{いつか}}
\textbf{\ja{いつか}} — \emph{itsuka} — \textquotedblleft{}the fifth of the month; some day\textquotedblright{}.

\begin{grammarlens}[title={What you've built}]
One more word for this chapter.
\end{grammarlens}

\subsection*{Your turn}
\begin{itemize}
\raggedright
  \item \emph{Say it:} \emph{itsuka}
  \item \emph{Say it:} \emph{itsuka}, once more
  \item \emph{Say it:} say \emph{itsuka}
  \item \emph{Recall:} say the Japanese for the third of the month, then the Japanese for the fourth of the month, then say \emph{itsuka} again
\end{itemize}

\begin{tcolorbox}[breakable,colback=teal!4,colframe=teal!35!black,title={Before you move on}]
What does \ja{いつか} mean? (\textquotedblleft{}The fifth of the month; some day\textquotedblright{}.) Say it once more.
\end{tcolorbox}
\section[\texorpdfstring{\ja{むいか} (muika)}{むいか (muika)}]{\ja{むいか} (muika) — the sixth of the month}

\label{lesson:JA-C77-muika}



\begin{quote}
Before the new one: say the Japanese for the fourth of the month, then the Japanese for the fifth of the month; some day.
\end{quote}

\subsection*{You'll want to know: \ja{むいか}}
\textbf{\ja{むいか}} — \emph{muika} — \textquotedblleft{}the sixth of the month\textquotedblright{}.

\begin{grammarlens}[title={What you've built}]
One more word for this chapter.
\end{grammarlens}

\subsection*{Your turn}
\begin{itemize}
\raggedright
  \item \emph{Say it:} \emph{muika}
  \item \emph{Say it:} \emph{muika}, once more
  \item \emph{Say it:} say \emph{muika}
  \item \emph{Recall:} say the Japanese for the fourth of the month, then the Japanese for the fifth of the month; some day, then say \emph{muika} again
\end{itemize}

\begin{tcolorbox}[breakable,colback=teal!4,colframe=teal!35!black,title={Before you move on}]
What does \ja{むいか} mean? (\textquotedblleft{}The sixth of the month\textquotedblright{}.) Say it once more.
\end{tcolorbox}
